\par\addvspace{12pt}\Needspace{9\baselineskip}

\section{单变量函数的连续性}

\label{sol:sec:functions:2}

% ITEM: functions-one-sided-piecewise
\solutionitem{functions-one-sided-piecewise}{例2.1}
设
\begin{equation*}
 f(x)=\begin{cases}
 x+1,&x<0,\\
 a,&x=0,\\
 x^2+b,&x>0.
 \end{cases}
\end{equation*}
求\(f\)在\(x=0\)处的左右极限，并分别求它在该点左连续、右连续及连续的条件。改变\(a\)能否改变左右极限？
\solutionanswer
左侧使用\(x+1\)，右侧使用\(x^2+b\)，所以左右极限分别为\(1\)、\(b\)。左连续的条件为\(a=1\)，右连续的条件为\(a=b\)，连续的条件为\(a=b=1\)。

改变\(a\)只改变零点处的取值，不改变任何去心邻域内的对应关系，因此不能改变左右极限。特别是\(b\ne1\)时，两侧极限不同，无论怎样改变\(a\)都不能使函数在零点连续。


% ITEM: functions-continuity-square
\solutionitem{functions-continuity-square}{习题2.1（1）}
用定义证明\(f(x)=x^2\)在\(x=1\)处连续。可先限制\(|x-1|<1\)，再控制\(|x+1|\)。
\solutionanswer
当\(|x-1|<1\)时，有\(|x+1|<3\)，故\(|x^2-1|\le3|x-1|\)。给定\(\varepsilon>0\)，取\(\delta=\min\{1,\varepsilon/3\}\)，则\(|x-1|<\delta\)时有\(|f(x)-f(1)|<\varepsilon\)。这正是连续性的定义。


% ITEM: functions-removable-hole
\solutionitem{functions-removable-hole}{习题2.1（2）}
当\(x\ne1\)时，设\(g(x)=(x^2-1)/(x-1)\)。应怎样定义\(g(1)\)，才能使它在\(1\)处连续？
\solutionanswer
当\(x\ne1\)时，\(g(x)=x+1\)，故\(\lim_{x\to1}g(x)=2\)。应定义\(g(1)=2\)，此时极限等于函数值，函数连续。


% ITEM: functions-bounded-factor-continuity
\solutionitem{functions-bounded-factor-continuity}{例2.2}
设\(g\)在\(x=0\)附近有界。定义\(f(0)=0\)，\(x\ne0\)时令\(f(x)=xg(x)\)。用连续性定义证明\(f\)在\(x=0\)处连续。是否需要\(g\)本身在\(x=0\)处连续？
\solutionanswer
取\(r>0,M>0\)，使\(0<|x|<r\)时有\(|g(x)|\le M\)。给定\(\varepsilon>0\)，取\(\delta=\min\{r,\varepsilon/M\}\)，则当\(0<|x|<\delta\)时，
\[
|f(x)-f(0)|=|xg(x)|\le M|x|<\varepsilon.
\]
\(x=0\)时误差为零，故\(f\)在零点连续。不需要\(g\)连续：例如令\(g(x)=1\)（\(x>0\)）、\(g(x)=-1\)（\(x<0\)）、\(g(0)=0\)，则\(g\)不连续，所得\(f(x)=|x|\)却连续。


% ITEM: functions-cubic-bisection
\solutionitem{functions-cubic-bisection}{例2.3}
证明方程\(x^3+x-1=0\)在\((0,1)\)内有唯一实根。不使用导数，怎样通过两点函数值之差判断严格单调？再写出进行三次二分后保留的区间。
\solutionanswer
令\(F(x)=x^3+x-1\)，则\(F(0)=-1\)、\(F(1)=1\)，由连续性和零点存在定理，区间内至少有一根。对任意\(u<v\)，
\[
F(v)-F(u)=(v-u)(v^2+uv+u^2+1)>0,
\]
因为\(v^2+uv+u^2=(v+u/2)^2+3u^2/4\ge0\)。所以\(F\)在\(\mathbb R\)上严格递增，实根唯一。

依次二分，\(F(1/2)=-3/8<0\)、\(F(3/4)=11/64>0\)、\(F(5/8)=-67/512<0\)，故三次后保留的区间为\([5/8,3/4]\)。


% ITEM: functions-fixed-point-without-convergence
\solutionitem{functions-fixed-point-without-convergence}{例2.4}
设\(f\)在\([0,1]\)上连续，且对所有\(x\in[0,1]\)都有\(0\le f(x)\le1\)。证明存在\(\xi\in[0,1]\)，使\(f(\xi)=\xi\)。这个结论是否保证从任意初值出发的迭代都收敛？
\solutionanswer
把目标写成零点问题，令\(h(x)=f(x)-x\)。它连续，且\(h(0)=f(0)\ge0\)、\(h(1)=f(1)-1\le0\)。若有端点为零，结论已成立；否则两端异号，由零点存在定理得到\(h(\xi)=0\)，即\(f(\xi)=\xi\)。

这不保证任意初值的迭代都收敛。例如\(f(x)=1-x\)满足条件，取初值\(0\)便得到\(0,1,0,1,\ldots\)。不动点的存在与迭代是否趋向它是不同的问题。


% ITEM: functions-cubic-onto
\solutionitem{functions-cubic-onto}{习题2.2（1）}
证明对任意实数\(t\)，方程\(x^3+x=t\)都有唯一实根。端点不再固定时，应怎样根据\(t\)选择一个有根的区间？
\solutionanswer
令\(F_t(x)=x^3+x-t\)。由上一例的两点差，\(F_t\)严格递增。取\(R=|t|+1\)，则\(F_t(-R)<0<F_t(R)\)。因此\([-R,R]\)内存在零点，而且唯一。区间随右端\(t\)改变，端点的选择只须使函数值超过\(|t|\)，不必先知道根的位置。


% ITEM: functions-graph-crossing
\solutionitem{functions-graph-crossing}{习题2.2（2）}
设\(f,g\)在\([a,b]\)上连续，且\(f(a)<g(a)\)、\(f(b)>g(b)\)。证明两函数的图像在\((a,b)\)内相交。
\solutionanswer
令\(h=f-g\)，则\(h\)在\([a,b]\)上连续，且\(h(a)<0<h(b)\)。由零点存在定理，存在\(\xi\in(a,b)\)使\(h(\xi)=0\)，即两图像在\((\xi,f(\xi))\)相交。


% ITEM: functions-zero-counterexamples
\solutionitem{functions-zero-counterexamples}{习题2.2（3）}
端点同号的连续函数能否在区间内部有零点？若没有连续性，端点异号是否仍能保证存在零点？分别举例说明。
\solutionanswer
端点同号也可能有零点。例如\(f(x)=(x-1/2)^2\)在\([0,1]\)两端都为正，在\(1/2\)处为零。

没有连续性，异号也不能保证有零点。例如在\([0,1]\)上令\(g(x)=-1\)（\(x<1/2\)）、\(g(x)=1\)（\(x\ge1/2\)），两端异号而处处不为零。


% ITEM: functions-quadratic-interlacing
\solutionitem{functions-quadratic-interlacing}{习题2.2（4）}
（2013年重庆卷理科第6题改编）若\(a<b<c\)，函数\(f(x)=(x-a)(x-b)+(x-b)(x-c)+(x-c)(x-a)\)的两个零点分别位于哪两个相邻区间内？
\solutionanswer
代入三个分界点，比直接求根更容易确定位置：
\[
f(a)=(a-b)(a-c)>0,\quad
f(b)=(b-c)(b-a)<0,\quad
f(c)=(c-a)(c-b)>0.
\]
故\((a,b)\)、\((b,c)\)内各有零点。\(f\)是二次函数，至多有两个零点，因此两个相邻区间内各恰有一个。


% ITEM: functions-open-closed-ranges
\solutionitem{functions-open-closed-ranges}{例2.5}
比较\(f(x)=x\)在\((0,1)\)与\([0,1]\)上的取值，再考察\(g(x)=1/x\)在\((0,1]\)上的取值。它们分别是否有界，是否有最大值和最小值？
\solutionanswer
\(f(x)=x\)在\((0,1)\)上的值域为\((0,1)\)，有界，但没有最大值和最小值；在\([0,1]\)上的值域为\([0,1]\)，最小值\(0\)在\(0\)处取得，最大值\(1\)在\(1\)处取得。\(g(x)=1/x\)在\((0,1]\)上的值域为\([1,+\infty)\)，向上无界，没有最大值，最小值为\(1\)。区间是否包含端点，直接影响边界值能否取得。


% ITEM: functions-quadratic-quantifiers
\solutionitem{functions-quadratic-quantifiers}{例2.6}
设\(f(x)=x^2-2x+2\)，\(x\in[0,3]\)。用配方法求值域，并分别求实数\(t\)的范围，使下列条件成立：

（1）对所有\(x\in[0,3]\)，都有\(f(x)\ge t\)。

（2）存在\(x\in[0,3]\)，使\(f(x)<t\)。

（3）方程\(f(x)=t\)在\([0,3]\)上有解，并进一步讨论不同\(t\)对应的根的个数。
\solutionanswer
配方得\(f(x)=(x-1)^2+1\)。在\([0,3]\)上，最小值为\(1\)，最大值为\(5\)，再由连续性，值域为\([1,5]\)。

（1）要使全部函数值都不小于\(t\)，须且只须\(t\le1\)。

（2）最小值能够取到，所以存在\(f(x)<t\)当且仅当\(t>1\)。

（3）有解当且仅当\(t\in[1,5]\)。由\(x=1\pm\sqrt{t-1}\)及原区间限制，根数为
\[
\begin{array}{c|c}
\text{参数范围}&\text{根数}\\\hline
 t<1\text{ 或 }t>5&0\\
 t=1&1\\
 1<t\le2&2\\
 2<t\le5&1
\end{array}
\]
减号对应的根在\(t>2\)时越出左端点，因而不能只看判别式来确定区间内的根数。


% ITEM: functions-parameter-nonnegative
\solutionitem{functions-parameter-nonnegative}{例2.7}
求实数\(t\)的范围，使不等式\(x^2-2tx+2\ge0\)对所有\(x\in[0,2]\)成立。再求存在\(x\in[0,2]\)使该表达式小于零的\(t\)的范围。将参数分离与直接求最小值两种方法进行比较。
\solutionanswer
先分离参数。\(x=0\)时原式为\(2>0\)；对\(0<x\le2\)，要求等价于
\[
t\le\frac{x^2+2}{2x}=\frac x2+\frac1x.
\]
由均值不等式，右边不小于\(\sqrt2\)，并在区间内的\(x=\sqrt2\)处取到。因此恒成立的参数范围为\(t\le\sqrt2\)。存在一个\(x\)使原式小于零的条件恰为\(t>\sqrt2\)，这时可直接取\(x=\sqrt2\)。

若直接求最小值，配方得到\((x-t)^2+2-t^2\)，还需考虑顶点\(t\)与区间的位置：
\[
\min_{[0,2]}(x^2-2tx+2)=
\begin{cases}2,&t\le0,\\2-t^2,&0\le t\le2,\\6-4t,&t\ge2.\end{cases}
\]
逐段判断同样得到\(t\le\sqrt2\)。本题中参数可以分离，先研究一个与\(t\)无关的表达式，便省去了顶点位置的分类。


% ITEM: functions-strict-quantifier-range
\solutionitem{functions-strict-quantifier-range}{习题2.3（1）}
求\(h(x)=x^2-4x+5\)在\([0,3]\)上的值域，并分别求使\(h(x)>t\)恒成立、至少存在一个\(x\)使\(h(x)>t\)的参数范围。
\solutionanswer
\(h(x)=(x-2)^2+1\)，在\([0,3]\)上最小值为\(1\)，最大值为\(5\)，值域为\([1,5]\)。由于最值均可取到，\(h(x)>t\)恒成立当且仅当\(t<1\)，至少有一个\(x\)满足\(h(x)>t\)当且仅当\(t<5\)。


% ITEM: functions-positive-compact-minimum
\solutionitem{functions-positive-compact-minimum}{习题2.3（2）}
设\(f\)在\([a,b]\)上连续且处处为正。证明存在常数\(c>0\)，使所有函数值都不小于\(c\)。将闭区间换为开区间，结论是否仍成立？
\solutionanswer
由最值定理，存在\(x_*\in[a,b]\)使\(f(x_*)\)为最小值。题设给出\(f(x_*)>0\)，故取\(c=f(x_*)\)即可。

换为开区间不一定成立。例如\(f(x)=x\)在\((0,1)\)上连续且处处为正，却能取任意接近零的正值，没有统一的正下界。


% ITEM: functions-equal-half-shift
\solutionitem{functions-equal-half-shift}{习题2.3（3）}
设\(f\)在\([0,1]\)上连续，且\(f(0)=f(1)\)。证明存在\(\xi\in[0,1/2]\)，使\(f(\xi+1/2)=f(\xi)\)。应对哪个函数使用零点存在定理？
\solutionanswer
目标是两处函数值之差为零，因此在\([0,1/2]\)上取\(h(x)=f(x+1/2)-f(x)\)。它连续，且
\[
h(0)=f(1/2)-f(0),\qquad h(1/2)=f(1)-f(1/2)=-h(0).
\]
若端点值为零，直接得到所求点；否则两端异号，由零点存在定理得到\(h(\xi)=0\)。


% ITEM: functions-alternating-quantifiers
\solutionitem{functions-alternating-quantifiers}{例2.8}
设连续函数\(f,g\)在两个闭区间\(I,J\)上的值域分别为\([m,M]\)、\([p,P]\)。分别求实数\(t\)的范围，使下列条件成立：

（1）对所有\(x\in I\)、\(y\in J\)，都有\(f(x)\ge t+g(y)\)。

（2）对每个\(x\in I\)，都存在\(y\in J\)，使\(f(x)\ge t+g(y)\)。

（3）存在\(x\in I\)，使对所有\(y\in J\)，都有\(f(x)\ge t+g(y)\)。
\solutionanswer
（1）左边可能最小为\(m\)，右边的\(g(y)\)可能最大为\(P\)，所以条件等价于\(m\ge t+P\)，即\(t\le m-P\)。必要性由最小值点和最大值点同时代入得到；充分性由\(f(x)\ge m\)、\(g(y)\le P\)得到。

（2）对任意给定的\(x\)，选择\(g\)的最小值点最有利。条件等价于所有\(f(x)\)都不小于\(t+p\)，即\(t\le m-p\)。若有某点满足\(f(x)<t+p\)，任何\(y\)都不能补救。

（3）要找一个\(x\)同时应对所有\(y\)，选择\(f\)的最大值点最有利；最大的右侧值为\(t+P\)。因此条件等价于\(M\ge t+P\)，即\(t\le M-P\)。


% ITEM: functions-exam-global-quadratic
\solutionitem{functions-exam-global-quadratic}{习题2.4（1）}
（2012年福建卷文科第15题）已知关于\(x\)的不等式\(x^2-ax+2a>0\)在\(\mathbb R\)上恒成立，则实数\(a\)的取值范围是\underline{\hspace{1.5cm}}。
\solutionanswer
配方得\(x^2-ax+2a=(x-a/2)^2+2a-a^2/4\)。它在\(\mathbb R\)上的最小值为\(a(8-a)/4\)。严格大于零恒成立当且仅当\(a(8-a)>0\)，故\(a\in(0,8)\)。


% ITEM: functions-exam-tianjin-parameter
\solutionitem{functions-exam-tianjin-parameter}{习题2.4（2）}
（2025年天津卷第15题）已知\(a,b\in\mathbb R\)，且\((2a+b)x^2+bx-a-1\le0\)对任意\(x\in[-2,2]\)成立，则\(2a+b\)的最小值为\underline{\hspace{1.2cm}}。
\solutionanswer
目标只含\(c=2a+b\)，应寻找能使剩余参数消去的自变量。用\(a=(c-b)/2\)改写原式，得到
\[
c\left(x^2-\frac12\right)+b\left(x+\frac12\right)-1\le0.
\]
取允许区间内的\(x=-1/2\)，\(b\)的系数消失，故\(-c/4-1\le0\)，即\(c\ge-4\)。

取\(a=0,b=-4\)，原式为\(-4x^2-4x-1=-(2x+1)^2\le0\)，对全部实数都成立。因此下界确可取到，所求最小值为\(\boxed{-4}\)。


% ITEM: functions-continuous-cubic-inverse
\solutionitem{functions-continuous-cubic-inverse}{例2.9}
证明\(f(x)=x^3+x\)在\(\mathbb R\)上有连续反函数，求其反函数的定义域和值域。若不求出反函数的表达式，能否判断\(\lim_{y\to2}f^{-1}(y)\)？
\solutionanswer
\(f\)是连续的多项式函数，且两点差证明它严格递增。对任意实数\(y\)，方程\(x^3+x=y\)都有唯一实根，故其值域为\(\mathbb R\)。由连续反函数定理，\(f^{-1}\)在\(\mathbb R\)上连续，定义域与值域均为\(\mathbb R\)。因为\(f(1)=2\)，所以
\[
\lim_{y\to2}f^{-1}(y)=f^{-1}(2)=1.
\]
这里只需确定原函数中与\(2\)对应的点，不需要解出一般三次方程。


% ITEM: functions-inverse-square-branches
\solutionitem{functions-inverse-square-branches}{习题2.5（1）}
将\(f(x)=x^2\)分别限制在\([-2,0]\)与\([0,2]\)上，写出两种反函数，并说明它们的单调方向与连续性。
\solutionanswer
在\([-2,0]\)上，原函数严格递减，值域为\([0,4]\)，反函数为\(f^{-1}(y)=-\sqrt y\)，定义域\([0,4]\)、值域\([-2,0]\)，严格递减且连续。

在\([0,2]\)上，反函数为\(f^{-1}(y)=\sqrt y\)，定义域\([0,4]\)、值域\([0,2]\)，严格递增且连续。两个分支对应不同的原函数定义域。


% ITEM: functions-monotone-disconnected-range
\solutionitem{functions-monotone-disconnected-range}{习题2.5（2）}
构造一个区间上的严格单调函数，使其值域不是区间。它的图像怎样说明，严格单调本身不能代替连续性？
\solutionanswer
在\(\mathbb R\)上取
\[
f(x)=\begin{cases}x,&x<0,\\x+1,&x\ge0.\end{cases}
\]
每一侧严格递增，跨过零点时函数值也严格增加，因此全域严格递增。但其值域为\(( -\infty,0)\cup[1,+\infty)\)，不是区间。图像在零点跳过了一段高度，说明严格单调只能保证一一对应，不能保证连续。


% ITEM: functions-inverse-rational-value
\solutionitem{functions-inverse-rational-value}{习题2.6（1）}
（2015年上海卷文科第4题）设\(f^{-1}(x)\)为\(f(x)=x/(2x+1)\)的反函数，则\(f^{-1}(2)=\underline{\hspace{1.2cm}}\)。
\solutionanswer
求\(f^{-1}(2)\)就是求\(x/(2x+1)=2\)的解。解得\(x=-2/3\)，分母为\(-1/3\ne0\)，所以\(\boxed{f^{-1}(2)=-2/3}\)。


% ITEM: functions-inverse-image-fixed-point
\solutionitem{functions-inverse-image-fixed-point}{习题2.6（2）}
（2013年上海卷理科第14题）对区间\(I\)上有定义的函数\(g(x)\)，记\(g(I)=\{y\mid y=g(x),x\in I\}\)。已知定义域为\([0,3]\)的函数\(y=f(x)\)有反函数\(y=f^{-1}(x)\)，且\(f^{-1}([0,1))=[1,2)\)，\(f^{-1}((2,4])=[0,1)\)。若方程\(f(x)-x=0\)有解\(x_0\)，则\(x_0=\underline{\hspace{1.2cm}}\)。
\solutionanswer
反函数的集合像条件说明：\(x\in[1,2)\)时\(f(x)\in[0,1)\)，\(x\in[0,1)\)时\(f(x)\in(2,4]\)。所以在前两个半开区间内均不可能有\(f(x)=x\)。

若\(x\in(2,3]\)，则\(x\in(2,4]\)。若又有\(f(x)=x\)，便由\(f^{-1}((2,4])=[0,1)\)推出\(x\in[0,1)\)，矛盾。定义域中只剩\(x=2\)未被排除，题设又保证存在解，因此\(\boxed{x_0=2}\)。这里没有使用连续性假设。


% ITEM: functions-oscillating-continuity
\solutionitem{functions-oscillating-continuity}{例2.10 *}
设\(f(x)=\sin(1/x)\)，\(x\ne0\)。分别选取趋于零的数列，使函数值恒等于\(1\)与\(-1\)，由此说明无论怎样定义\(f(0)\)，都不能使它在\(x=0\)处连续。再比较\(g(x)=x\sin(1/x)\)、\(g(0)=0\)的情形。
\solutionanswer
取\(x_n=(\pi/2+2n\pi)^{-1}\)、\(y_n=(3\pi/2+2n\pi)^{-1}\)。两列均趋于零，而\(f(x_n)=1\)、\(f(y_n)=-1\)。若零点处能够连续，两列函数值应趋于同一个\(f(0)\)，矛盾。

对\(g\)，有\(|g(x)|\le|x|\)，故\(x\to0\)时\(g(x)\to0=g(0)\)。原来的振荡仍然存在，但振幅随\(|x|\)趋于零，因而不破坏连续性。


% ITEM: functions-iteration-fixed-point
\solutionitem{functions-iteration-fixed-point}{习题2.7（1） *}
设\(f\)在\([0,1]\)上连续，数列\(x_n\)满足\(x_{n+1}=f(x_n)\)，且各项都在\([0,1]\)内。
若\(x_n\)收敛，证明其极限必为\(f\)的不动点。
\solutionanswer
设\(x_n\to L\)。因各项在\([0,1]\)内，\(L\in[0,1]\)；由连续性，在\(x_{n+1}=f(x_n)\)中取极限，得\(L=f(L)\)。这先以收敛为前提，确定极限必须满足的方程。


% ITEM: functions-reflection-iteration
\solutionitem{functions-reflection-iteration}{习题2.7（2） *}
设\(f\)在\([0,1]\)上连续，数列\(x_n\)满足\(x_{n+1}=f(x_n)\)，且各项都在\([0,1]\)内。
取\(f(x)=1-x\)，讨论不同初值产生的数列，并说明不动点的存在为什么不能保证迭代收敛。
\solutionanswer
设初值为\(c\in[0,1]\)，则奇数项均为\(c\)，偶数项均为\(1-c\)。当\(c=1/2\)时为常数列，收敛于\(1/2\)；其他初值产生两点交替的数列，不收敛。不动点存在只说明有一个保持不动的初值，并未保证其他初值趋向它。


% ITEM: functions-two-cycle-limits
\solutionitem{functions-two-cycle-limits}{习题2.7（3） *}
设\(f\)在\([0,1]\)上连续，数列\(x_n\)满足\(x_{n+1}=f(x_n)\)，且各项都在\([0,1]\)内。
若偶数项与奇数项分别趋于\(u,v\)，证明\(f(u)=v\)、\(f(v)=u\)。此时整个数列收敛的充要条件是什么？结合上一问解释\(u\ne v\)时的行为。
\solutionanswer
由\(x_{2n+1}=f(x_{2n})\)、\(x_{2n+2}=f(x_{2n+1})\)和连续性，取极限得\(v=f(u)\)、\(u=f(v)\)。原数列收敛当且仅当\(u=v\)：必要性来自子列同极限，充分性来自奇偶两列合并的收敛判据。当\(u\ne v\)时，两点由\(f\)互相映到对方；上一问的\(c,1-c\)正给出这样的两点循环。


% ITEM: functions-uniform-sqrt
\solutionitem{functions-uniform-sqrt}{习题2.8（1） *}
证明\(f(x)=\sqrt x\)在\([0,+\infty)\)上一致连续。可先证明\(|\sqrt x-\sqrt y|\le\sqrt{|x-y|}\)。
\solutionanswer
设\(x\ge y\ge0\)，则
\[
(\sqrt x-\sqrt y)^2\le(\sqrt x-\sqrt y)(\sqrt x+\sqrt y)=x-y.
\]
交换\(x,y\)后同样成立，因此\(|\sqrt x-\sqrt y|\le\sqrt{|x-y|}\)。给定\(\varepsilon>0\)，统一取\(\delta=\varepsilon^2\)，即可使全部满足\(|x-y|<\delta\)的点对都有函数值之差小于\(\varepsilon\)。


% ITEM: functions-uniform-reciprocal
\solutionitem{functions-uniform-reciprocal}{习题2.8（2） *}
判断\(g(x)=1/x\)在\([1,+\infty)\)与\((0,1)\)上是否一致连续，并给出证明。
\solutionanswer
在\([1,+\infty)\)上，
\[
\left|\frac1x-\frac1y\right|=\frac{|x-y|}{xy}\le|x-y|,
\]
故取\(\delta=\varepsilon\)证明一致连续。

在\((0,1)\)上，取\(x_n=1/(n+1)\)、\(y_n=1/(n+2)\)。两点距离趋于零，函数值之差却恒为\(1\)，由一致连续的数列刻画可知不一致连续。


% ITEM: functions-uniform-cauchy-image
\solutionitem{functions-uniform-cauchy-image}{习题2.8（3） *}
设\(f\)在\(A\)上一致连续，\(x_n\)是\(A\)中的Cauchy数列。证明\(f(x_n)\)也是Cauchy数列。用\(g(x)=1/x\)及\(x_n=1/(n+1)\)说明，逐点连续不能保证这一结论。
\solutionanswer
给定\(\varepsilon>0\)，由一致连续取\(\delta>0\)。再由\(x_n\)为Cauchy数列，取\(N\)使\(m,n>N\)时\(|x_m-x_n|<\delta\)。于是\(|f(x_m)-f(x_n)|<\varepsilon\)，所以函数值也构成Cauchy数列。

\(g(x)=1/x\)在\((0,1)\)上逐点连续，而\(x_n=1/(n+1)\)是这个区间内的Cauchy数列，\(g(x_n)=n+1\)却不是Cauchy数列。逐点连续在每个基点附近提供的范围不同，不能代替上述对全部点对的统一控制。
\par
